% =====================================================================
%  第 3 章  智能体设计
% =====================================================================

\section{智能体设计}
\label{sec:method}

本章描述系统的设计方案。系统骨架尚未完全实现，因此本章不含实验数据，部分实现细节将在开发过程中调整。

设计之前须先回答一个问题：本任务为何需要智能体，而不能由固定脚本完成。一个固定流程的程序亦可完成本任务：枚举不等式形式，以线性规划确定系数，在数据上筛选并输出结果。该流程不需要任何智能，但其效果受限，原因在于流程中存在若干无法预先固定的环节。

第一，搜索方向的确定。图的不变量有十余个，两两组合即构成上百种形式，穷举既缓慢又多数无意义。有效做法是先考察数据中哪些量之间存在异常关联，再针对性地提出形式。

第二，反例搜索失败的解释。未找到反例可能意味着候选成立，也可能意味着搜索策略不当，二者的区分依赖于对失败原因的理解。若始终在随机图上搜索，而反例只可能出现于特定结构中，则“未找到”不提供任何信息。

第三，候选新颖性的判定。部分“新发现”实为文献中已有定理的改写。判定这一点需要语义层面的理解，数值比对无法胜任。

上述环节都需要在开放空间中作出判断，规则无法穷举，一旦穷举即退化为固定脚本。因此系统由大语言模型承担此类决策，由确定性代码承担计算与验证。系统的形态并非由大语言模型全程生成，而是由其在关键节点作出决策、由代码执行与校验。

\subsection{总体结构}

系统以闭环方式运行，每一轮包含四个阶段，其结构如图 \ref{fig:loop} 所示。

\begin{figure}[H]
\centering
\begin{tikzpicture}[
  node distance=0.55cm,
  box/.style={draw=accent!55, fill=accentlight, rounded corners=2pt,
              minimum width=2.55cm, minimum height=0.85cm,
              align=center, font=\small},
  arr/.style={-{Stealth[length=2mm]}, draw=accent!70, thick},
  fb/.style={-{Stealth[length=2mm]}, draw=orange!75!black, thick, dashed}
]
\node[box] (a) {提出候选\\{\scriptsize 大模型与线性规划}};
\node[box, right=of a] (b) {搜索反例\\{\scriptsize 四类策略}};
\node[box, right=of b] (c) {统计筛选\\{\scriptsize 控制假发现率}};
\node[box, below=0.95cm of c] (d) {记录与反思\\{\scriptsize 写入轨迹，修正下一轮}};

\draw[arr] (a) -- (b);
\draw[arr] (b) -- (c);
\draw[arr] (c) -- (d);
\draw[fb] (d.west) -| node[pos=0.25, below, font=\scriptsize, text=orange!70!black] {失败原因反馈} (a.south);
\end{tikzpicture}
\caption{智能体的单轮运行结构}
\label{fig:loop}
\end{figure}

\textbf{提出候选。}将当前数据的概况（各不变量的取值范围、两两之间的相关性、已确认成立的定理清单）输入大语言模型，由其提出若干值得尝试的不等式形式。模型仅决定使用哪几个不变量及各自的方向，具体系数由线性规划求解，无需由模型给出。

\textbf{搜索反例。}对每条候选，在有限算力内搜索使其不成立的图。搜索分为四类：随机图、特殊结构图（如二分图与完全多部图）、由已知临界例子作局部修改、以及小规模图的完全枚举。各类策略的预算分配由上一轮的成败情况决定。

\textbf{统计筛选。}将“在多少张图上成立”转换为统计意义上的可信程度，并在同时检验大量候选时作多重检验校正，以控制最终保留集合中无效命题的比例。

\textbf{记录与反思。}每一步的输入、输出与决策均写入结构化日志。筛选阶段暴露的失败模式被送回第一阶段，用于修正下一轮的形式提议。

该反馈回路是系统区别于固定流水线的实质所在：系统的策略随运行结果而调整。

\subsection{大语言模型的职责}

系统中共有四处调用大语言模型，其职责如下。

\begin{enumerate}
  \item \textbf{提出形式}。输入为数据摘要与已知定理清单，输出为若干候选形式，以结构化格式描述所涉不变量、各自符号及是否含常数项。形式空间是开放的，无法以规则穷举，故此项为核心。
  \item \textbf{解释失败}。当候选被反例推翻时，将反例的结构特征输入模型，由其归纳该候选失效的原因。所得结论作为下一轮的上下文，以避免重复提出同类错误形式。
  \item \textbf{判定新颖性}。将候选形式转写为自然语言描述，由模型对照已知定理清单判断其是否为某个已知结论的改写。
  \item \textbf{选择搜索策略}。给定候选与上一轮搜索记录，由模型判断应重点搜索的方向。
\end{enumerate}

以下环节不交由大语言模型处理：系数求解、反例的实际搜索、统计检验的计算、结果的记录。这些环节均有确定算法，交由模型只会引入不确定性。

为检验大语言模型的实际贡献，拟设计消融实验：将第一阶段的模型调用替换为随机选择形式，其余环节保持不变，在相同预算下比较两者所产出候选的质量，指标包括通过筛选的比例与在被反例推翻前存活的轮数。若二者无显著差异，则说明模型在该任务中贡献有限，此结论将如实报告。

\subsection{工具层与运行轨迹}

系统通过一组工具与外部交互。这些工具均为普通函数，由代码实现，由大语言模型决定调用时机，其功能如表 \ref{tab:tools} 所示。

\begin{table}[H]
\centering
\small
\caption{工具层接口}
\label{tab:tools}
\begin{tabular}{@{}ll@{}}
\toprule
工具 & 功能 \\
\midrule
\texttt{graph\_sample} & 生成图样本（随机图、特殊结构图、小图枚举） \\
\texttt{invariants} & 计算一张图的各项不变量 \\
\texttt{lp\_solve} & 给定形式，以线性规划求解系数 \\
\texttt{counterexample} & 在给定预算内搜索反例 \\
\texttt{permutation\_test} & 构造零分布并计算 $p$ 值 \\
\texttt{bh\_filter} & 作多重检验校正，输出保留集合 \\
\texttt{trace\_write} & 将本轮记录写入 JSONL 日志 \\
\bottomrule
\end{tabular}
\end{table}

其中 \texttt{counterexample} 的计算量最大，也是唯一需要严格限时的环节。系统为其设定固定预算（例如每条候选 30 秒），超时即返回未找到，并同时记录已用预算。搜索 30 秒未找到与搜索 3 小时未找到所提供的证据强度不同，这一区别必须在数据中保留。

整个循环的每一步写入一个 JSONL 文件，每行一条记录，字段如表 \ref{tab:trace} 所示。

\begin{table}[H]
\centering
\small
\caption{运行轨迹的记录字段}
\label{tab:trace}
\begin{tabular}{@{}ll@{}}
\toprule
字段 & 含义 \\
\midrule
\texttt{round} & 轮次 \\
\texttt{stage} & 所处阶段（提议、反驳、筛选） \\
\texttt{candidate} & 候选的形式与系数 \\
\texttt{llm\_input} / \texttt{llm\_output} & 模型的输入提示与原始输出 \\
\texttt{strategy} / \texttt{budget} & 本次搜索所用策略与预算 \\
\texttt{counterexample} & 找到的反例（若有） \\
\texttt{p\_value} / \texttt{decision} & 统计筛选结果与最终判定 \\
\bottomrule
\end{tabular}
\end{table}

该日志有三项用途：复现，即依据日志重放全部决策过程；分析，即考察各类策略的有效性及失败集中的形式类别；汇报，即由日志直接生成图表，无需人工整理。

需要说明的是，本系统包含大量固定流程，大语言模型仅占据上述四个决策点。若将这四点全部替换为随机选择，系统仍可运行，只是效果下降。因此本系统不宜称为完全自主的智能体，其准确定位是以大语言模型为决策核心、以确定性算法为执行手段的混合系统。
